%! TEX program = lualatex
\documentclass[14pt]{extarticle}

\usepackage[ukrainian]{babel}
\usepackage[margin=2cm]{geometry}

\usepackage{fontspec}
\setmainfont{CMU Serif}
\setsansfont{CMU Sans Serif}
\defaultfontfeatures{Ligatures={TeX}}

\usepackage[hyphens]{url}

\addto\captionsukrainian{\renewcommand{\refname}{Список використаних джерел}}

\emergencystretch=1em
\hbadness=10000
\pagestyle{empty}

\begin{document}

\section*{Огляд літератури (bibtex8, стиль plain)}

Загальні принципи адитивних технологій описано у~\cite{Gibson2021}.
Проектування FFF-принтерів розглянуто у~\cite{Soni2021},
а порівняльний аналіз картезіанських та дельта-принтерів~--- у~\cite{Schmitt2018}.
Типові дефекти FDM-друку систематизовано у~\cite{Baechle2022}.
Методи виявлення помилок на основі комп'ютерного зору досліджено
у~\cite{Baumann2016, Langeland2020, Paulsen2019}.
Машинне навчання для контролю якості поверхні застосовано у~\cite{Kadam2021}.
Методи оцінки якості зображень описано у~\cite{Crete2007, Wang2005}.
Основи цифрової обробки зображень викладено у~\cite{Gonzalez2018}.
Технічні характеристики комерційних 3D-принтерів наведено у~\cite{BambuLabA1Mini}.
Моделювання вбудованого 3D-друку розглянуто у~\cite{VanDriel2025}.

\nocite{*}

\bibliographystyle{plain}
\bibliography{references}

\end{document}
